\documentclass[10pt,a4paper]{article}

\usepackage[a4paper,margin=18mm,headheight=14pt]{geometry}
\usepackage{array}
\usepackage{booktabs}
\usepackage{enumitem}
\usepackage{fancyhdr}
\usepackage{float}
\usepackage{fontspec}
\usepackage{graphicx}
\usepackage{longtable}
\usepackage{microtype}
\usepackage{parskip}
\usepackage{xcolor}
\usepackage{hyperref}

\definecolor{ink}{HTML}{1F2937}
\definecolor{muted}{HTML}{64748B}
\definecolor{accent}{HTML}{0F766E}
\definecolor{panel}{HTML}{F1F5F9}
\definecolor{pending}{HTML}{FEF3C7}

\hypersetup{
  colorlinks=true,
  linkcolor=accent,
  urlcolor=accent,
  pdftitle={TokTickIT Lab 1 Report},
  pdfauthor={TokTickIT},
}

\setmainfont{Latin Modern Roman}
\setlength{\parindent}{0pt}
\setlist{nosep,leftmargin=*}
\setlist[itemize]{label=\textcolor{accent}{\small\textbullet}}

\pagestyle{fancy}
\fancyhf{}
\lhead{TokTickIT Lab 1}
\rhead{Submission evidence}
\cfoot{\thepage}
\renewcommand{\headrulewidth}{0.3pt}

\newcommand{\captureplaceholder}[1]{%
  \fcolorbox{muted}{pending}{%
    \parbox[c][0.22\textheight][c]{0.94\linewidth}{%
      \centering\textcolor{ink}{\textbf{Evidence pending}}\\[3pt]
      \small #1
    }%
  }%
}

\newcommand{\evidenceimage}[2]{%
  \IfFileExists{#1}{%
    \includegraphics[width=\linewidth,height=0.22\textheight,keepaspectratio]{#1}%
  }{\captureplaceholder{#2}}%
}

\newcommand{\workflowimage}[2]{%
  \IfFileExists{#1}{%
    \includegraphics[width=\linewidth,height=0.30\textheight,keepaspectratio]{#1}%
  }{\captureplaceholder{#2}}%
}

\newcommand{\identitycard}[5]{%
  \fcolorbox{accent}{panel}{%
    \parbox[t]{\dimexpr\linewidth-2\fboxsep-2\fboxrule\relax}{%
      {\small\bfseries\textcolor{accent}{#1}}\par
      {\bfseries #2}\par
      {\small ID: #3}\par
      {\small GitHub: \href{#4}{#5}}%
    }%
  }%
}

\newcommand{\promptquote}[1]{%
  \textcolor{accent}{\textquotedblleft}\emph{#1}\textcolor{accent}{\textquotedblright}%
}

\begin{document}

\begin{titlepage}
  \thispagestyle{empty}
  \vspace*{18mm}
  {\Large\textcolor{accent}{CPE 334 \textbar\ Individual Sprint 1}}\par
  \vspace{10mm}
  {\Huge\bfseries TokTickIT}\par
  \vspace{3mm}
  {\Large Lab 1 full-stack vertical slice}\par
  \vfill
  \begin{tabular}{@{}>{\bfseries}p{35mm}p{95mm}@{}}
    Stack & React, TypeScript, Vite, Bootstrap, Express, Prisma, PostgreSQL \\
    Workflow & Bun workspaces, feature branches, GitHub Issues, Pull Requests \\
    Evidence & Tests, API contracts, database seed, and human-checked App Demo \\
  \end{tabular}
  \vfill
  {\small Built from one LaTeX source with XeLaTeX.}
\end{titlepage}

\tableofcontents
\clearpage

\section{Answer Part 1: Git Use with Engineering Workflow}

The repository uses four parent Lab 1 issues and native sub-issues. Work is performed on feature branches and targets \texttt{lab1-staging}; the stable branch is \texttt{main}.

\begin{tabular}{@{}p{31mm}p{52mm}p{42mm}@{}}
  \toprule
  Parent issue & Feature branch & Pull Request \\
  \midrule
  \href{https://github.com/ovenmakemeheat/toktickit/issues/13}{\#13 foundation} & \href{https://github.com/ovenmakemeheat/toktickit/tree/feature/1-project-foundation}{\texttt{feature/1-project-foundation}} & \href{https://github.com/ovenmakemeheat/toktickit/pull/44}{PR \#44} \\
  \href{https://github.com/ovenmakemeheat/toktickit/issues/14}{\#14 health} & \href{https://github.com/ovenmakemeheat/toktickit/tree/feature/2-health-check}{\texttt{feature/2-health-check}} & \href{https://github.com/ovenmakemeheat/toktickit/pull/46}{PR \#46} \\
  \href{https://github.com/ovenmakemeheat/toktickit/issues/15}{\#15 categories} & \href{https://github.com/ovenmakemeheat/toktickit/tree/feature/3-category-seed}{\texttt{feature/3-category-seed}} & \href{https://github.com/ovenmakemeheat/toktickit/pull/45}{PR \#45} \\
  \href{https://github.com/ovenmakemeheat/toktickit/issues/16}{\#16 category list} & \href{https://github.com/ovenmakemeheat/toktickit/tree/feature/4-category-list}{\texttt{feature/4-category-list}} & \href{https://github.com/ovenmakemeheat/toktickit/pull/47}{PR \#47} \\
  \bottomrule
\end{tabular}

\vspace{4mm}
\begin{tabular}{@{}p{35mm}p{125mm}@{}}
  Repository & \href{https://github.com/ovenmakemeheat/toktickit}{ovenmakemeheat/toktickit} \\
  Project board & \href{https://github.com/users/ovenmakemeheat/projects/1/views/2}{TokTickIT Lab 1} \\
  Integration branch & \href{https://github.com/ovenmakemeheat/toktickit/tree/lab1-staging}{\texttt{lab1-staging}} \\
  Stable branch & \href{https://github.com/ovenmakemeheat/toktickit/tree/main}{\texttt{main}} \\
  Integration status & All four feature PRs are merged into \texttt{lab1-staging}; no staging-to-main PR exists in the audited GitHub state. \\
\end{tabular}

\vspace{3mm}
\textbf{Student and peer reviewer}

\vspace{1mm}
\begin{tabular}{@{}p{0.47\linewidth}@{\hspace{4mm}}p{0.47\linewidth}@{}}
  \identitycard{Student}{GUNTEE DOUNGMANEE}{67070501003}{https://github.com/ovenmakemeheat}{ovenmakemeheat}
  &
  \identitycard{Peer reviewer}{POLWARIT WATTHANAHEMMARAT}{67070501067}{https://github.com/MadMax168}{MadMax168}
\end{tabular}

\vspace{3mm}
\begin{itemize}
  \item Repository: \url{https://github.com/ovenmakemeheat/toktickit}
  \item Project board: \url{https://github.com/users/ovenmakemeheat/projects/1/views/2}
  \item The parent issue checklists are the acceptance source; native sub-issues carry implementation, testing, and evidence slices.
  \item Positive peer-review records for PRs \#44-\#47 are audited with the author responses.
  \item Full review details are in \texttt{docs/lab-01/reviewer.md}; PRs \#44-\#47 are merged into \texttt{lab1-staging}.
  \item Unresolved automated threads on PRs \#44, \#45, and \#47 are recorded separately and are not peer-review rejections.
  \item Reciprocal review target: \href{https://github.com/MadMax168/TikTockIT}{MadMax168/TikTockIT}; its public pull-request list was empty at audit time, so reciprocal review evidence remains a human action.
\end{itemize}

\textbf{Student--peer reviewer conversation}

{\small
\begin{longtable}{@{}>{\raggedright\arraybackslash}p{29mm}@{\hspace{2mm}}>{\raggedright\arraybackslash}p{0.34\linewidth}@{\hspace{2mm}}>{\raggedright\arraybackslash}p{0.42\linewidth}@{}}
  \toprule
  \textbf{Pull request} & \textbf{Peer reviewer message} & \textbf{Student response} \\
  \midrule
  \endfirsthead
  \toprule
  \textbf{Pull request} & \textbf{Peer reviewer message} & \textbf{Student response} \\
  \midrule
  \endhead
  \href{https://github.com/ovenmakemeheat/toktickit/pull/44}{\#44 foundation} \textit{(Merged)} & \promptquote{All things good for me} & Acknowledged the lockfile comment and stated that the synchronized lockfile would be regenerated and committed. \\
  \href{https://github.com/ovenmakemeheat/toktickit/pull/45}{\#45 category seed} \textit{(Merged)} & \promptquote{Good for me} & \promptquote{LGTM} \\
  \href{https://github.com/ovenmakemeheat/toktickit/pull/46}{\#46 health check} \textit{(Merged)} & \promptquote{Nothing to change, Good for Me} & \promptquote{Thanks for the review!} \\
  \href{https://github.com/ovenmakemeheat/toktickit/pull/47}{\#47 category list} \textit{(Merged)} & \promptquote{All Be Good for Me} & \promptquote{Thanks!} The pull request was merged into \texttt{lab1-staging}. \\
  \bottomrule
\end{longtable}
}

\emph{GitHub records these peer submissions as \texttt{COMMENTED}; they are conversation evidence, not formal \texttt{APPROVE} events.}

\textbf{Repository workflow screenshots}

The following captures provide human-checked evidence for the board state, merge history, repository structure, and automated test output. GitHub captures were taken with browser-use CLI from the relevant \texttt{lab1-staging} and \texttt{main} views.

\begin{figure}[H]
  \centering
  \workflowimage{docs/lab-01/evidence/board.png}{Capture the completed Lab 1 project board.}
  \caption{Lab 1 project board: six workflow columns, all four parent issues in Done, and 100\% sub-issue progress.}
\end{figure}

\begin{figure}[H]
  \centering
  \workflowimage{docs/lab-01/evidence/commit-history.png}{Capture the lab1-staging commit and merge history.}
  \caption{Commit and merge history on \texttt{lab1-staging}, including the PR \#44 merge and issue-referenced commits.}
\end{figure}

\begin{figure}[H]
  \centering
  \workflowimage{docs/lab-01/evidence/main-history.png}{Capture the current main branch history.}
  \caption{Current \texttt{main} history: it still contains only the initial setup commits, so staging-to-main integration remains a human action.}
\end{figure}

\begin{figure}[H]
  \centering
  \workflowimage{docs/lab-01/evidence/directory-structure.png}{Capture the repository directory structure.}
  \caption{Repository tree showing the client, server, documentation, agent instructions, and root tooling.}
\end{figure}

\begin{figure}[H]
  \centering
  \workflowimage{docs/lab-01/evidence/docs-lab1-structure.png}{Capture the docs/lab-01 directory contents.}
  \caption{Nested documentation tree showing \texttt{tests.md}, \texttt{reviewer.md}, \texttt{ai\_use.md}, and the LaTeX report source.}
\end{figure}

\begin{figure}[H]
  \centering
  \workflowimage{docs/lab-01/evidence/server-tests-structure.png}{Capture the server/tests/lab-01 directory contents.}
  \caption{Nested server test tree showing the Lab 1 API, seed, foundation, and setup test files.}
\end{figure}

\begin{figure}[H]
  \centering
  \workflowimage{docs/lab-01/evidence/test-capture.png}{Capture passing client and server test output.}
  \caption{Terminal evidence from the integrated branch: the client suite passes 3 tests and the server suite passes 6 tests.}
\end{figure}

\textbf{Rendered README.md}

\begin{tabular}{@{}p{31mm}p{129mm}@{}}
  Purpose & TokTickIT is the Lab 1 full-stack foundation for the IT service desk application. \\
  Stack & Bun workspaces; React, TypeScript, Vite, Bootstrap; Node-compatible Express and TypeScript; PostgreSQL 16 through Docker Compose; Prisma; Biome; Lefthook; Vitest; and Supertest. \\
  Prerequisites & Bun 1.3 or newer and Docker with Compose. \\
  Setup & \texttt{bun install}; copy \texttt{server/.env.example} to \texttt{server/.env}; install hooks; start the database; generate and validate Prisma. \\
  Database & The stack is owned by \texttt{server/}; Compose is under \texttt{server/docker-compose.yml}, PostgreSQL initialization is under \texttt{server/docker/}, and Prisma files are under \texttt{server/prisma/}. \\
  Seed and tests & \texttt{bun run db:migrate}, \texttt{bun run db:seed}, \texttt{bun run db:test:setup}, then \texttt{bun run verify}. \\
  Run & \texttt{bun run dev}; client at \url{http://localhost:5173}, API at \url{http://localhost:3000}, and PostgreSQL at \texttt{localhost:5432}. \\
  Scope & Lab 1 proves the client, API, Prisma, PostgreSQL, and test tooling as one vertical slice. Authentication, tickets, uploads, and later-lab screens are out of scope. \\
\end{tabular}

\vspace{2mm}
\textbf{Rendered .gitignore}

\begin{tabular}{@{}p{38mm}p{122mm}@{}}
  Dependencies & \texttt{node\_modules/}, \texttt{.bun/}, and package-manager stores. \\
  Environment & \texttt{.env} and \texttt{.env.*}, except \texttt{.env.example}. \\
  Build and cache & \texttt{dist/}, \texttt{build/}, \texttt{.vite/}, coverage, TypeScript build info, and logs. \\
  Local state & \texttt{*.db}, \texttt{*.db-journal}, \texttt{tmp/}, and OS metadata. \\
  Editor and agent files & Local VS Code settings, \texttt{.claude/}, and \texttt{CLAUDE.md}. \\
  LaTeX artifacts & \texttt{*.aux}, \texttt{*.fdb\_latexmk}, \texttt{*.fls}, \texttt{*.synctex.gz}, \texttt{*.toc}, \texttt{*.out}, and \texttt{*.xdv}. \\
\end{tabular}

\section{Answer Part 2: Tests}

The test matrix is kept in \texttt{docs/lab-01/tests.md}. The application boundary is tested through the exported Express app and user-observable React behavior.

\begin{tabular}{@{}p{13mm}p{42mm}p{16mm}p{55mm}p{14mm}@{}}
  \toprule
  ID & Test file & Tool & Proof & Status \\
  \midrule
  API-01 & \path{server/tests/lab-01/health.test.ts} & Supertest & Health returns HTTP 200 and the exact TokTickIT API payload. & Passing \\
  API-02 & \path{server/tests/lab-01/categories.test.ts} & Supertest & Categories return seeded ID/name pairs in ascending ID order. & Passing \\
  UI-01 & \path{client/src/App.test.tsx} & Vitest & TokTickIT heading and Check System action render. & Passing \\
  UI-02 & \path{client/src/App.test.tsx} & Vitest & Checking transitions to Online and renders returned categories. & Passing \\
  UI-03 & \path{client/src/App.test.tsx} & Vitest & API failure clears stale categories and shows the safe error. & Passing \\
  \bottomrule
\end{tabular}

\vspace{4mm}
The additional foundation smoke test is at \path{server/tests/lab-01/foundation.test.ts}; it verifies the exported Express app and typed environment configuration. This table is the rendered test traceability from \path{docs/lab-01/tests.md}.

\vspace{2mm}
\textbf{Verification command:} \texttt{bun run verify}. It covers Biome, Lefthook validation, Prisma validation, type checks, both workspace test suites, and both workspace builds.

\section{Answer Part 3: AI Use and Reflection}

\textbf{Tool and model}

\begin{tabular}{@{}p{34mm}p{124mm}@{}}
  \toprule
  Tool & OpenAI Codex coding agent in ChatGPT \\
  Model & GPT-5 \\
  Supporting tools & GitHub connector, browser-use CLI, Bun, XeLaTeX, and Poppler \\
  \bottomrule
\end{tabular}

\vspace{2mm}
\textbf{Selected prompts and effects}

{\small
\begin{longtable}{@{}>{\raggedright\arraybackslash}p{7mm}@{\hspace{2mm}}>{\raggedright\arraybackslash}p{0.37\linewidth}@{\hspace{2mm}}>{\raggedright\arraybackslash}p{0.52\linewidth}@{}}
  \toprule
  \textbf{ID} & \textbf{Selected prompt} & \textbf{Effect} \\
  \midrule
  \endfirsthead
  \toprule
  \textbf{ID} & \textbf{Selected prompt} & \textbf{Effect} \\
  \midrule
  \endhead
  P1 & \promptquote{read @docs/requirements, \$grill-with-docs} & Loaded the requirements and used the grilling workflow to identify scope and risks. \\
  P2 & \promptquote{now follow the plan, first create the issues} & Converted the requirements into the Lab 1 issue structure. \\
  P3 & \promptquote{all subissues should be inside the main 4 issues} & Kept exactly four parent issues and attached implementation work as native sub-issues. \\
  P4 & \promptquote{please keep insturction in AGENTS.md strictly about don't close any PR or main issue you can only do check the checklists or closing subissue, closing main issue / PR close is human job} & Established human-only rules for pull-request and main-issue state changes. \\
  P5 & \promptquote{add biome, left-hook, prisma validation on git sub-issues of issue 1} & Added the requested quality gates and connected them to the foundation work. \\
  P6 & \promptquote{yes, and pdf should be built in latex} & Chose a reproducible XeLaTeX report source and PDF build command. \\
  P7 & \promptquote{do capture those for me from http://localhost:5173/} & Defined the required loading, success, and failure evidence for the running app. \\
  P8 & \promptquote{do use browseruse cli} & Used browser-use CLI to verify the visible UI states and save the captures. \\
  P9 & \promptquote{this is student - reviewer details Student ... also checkout all pr and review records; update reviewer part} & Audited PRs \#44-\#47 and recorded student, peer, review comments, responses, and unresolved automated threads. \\
  \bottomrule
\end{longtable}
}

\textbf{Reflection}

The prompts became more useful as they added concrete constraints: Bun, the server-owned database stack, four parent issues, native sub-issues, acceptance criteria, and human-only pull-request state changes. Those constraints kept the implementation aligned with Lab 1 instead of expanding into later features.

The AI-assisted workflow still required human-in-the-loop checks. The browser capture was inspected visually; an initial loading attempt was detected as a completed success state and recaptured correctly. The final PDF was rendered page by page and checked for readable evidence, clipping, overlap, and excessive empty space. GitHub review records were also checked directly so positive peer comments were not mislabeled as formal approvals, and unresolved automated threads were preserved rather than reported as fixed.

\section{Answer Part 4: App Demo}

The system check makes two relative API requests through the Vite proxy: \path{/api/health} for health and \path{/api/categories} for categories. The success state proves that health is online and that the four categories came from PostgreSQL. The failure state proves that the user receives a safe error and stale categories are cleared.

\begin{figure}[H]
  \centering
  \evidenceimage{docs/lab-01/evidence/loading.png}{Capture the transient Checking state while requests are in flight.}
  \caption{Loading state: the button and status show that the system check is in progress.}
\end{figure}

\begin{figure}[H]
  \centering
  \evidenceimage{docs/lab-01/evidence/success.png}{Capture Online status and all four returned categories.}
  \caption{Success state: Online and the database-backed categories are visible in order.}
\end{figure}

\begin{figure}[H]
  \centering
  \evidenceimage{docs/lab-01/evidence/failure.png}{Capture Offline status, the safe error, and no stale category list.}
  \caption{Failure state: the API is unavailable and the user-safe error is shown.}
\end{figure}


\end{document}
